% GENERATED FILE. Edit canonical lessons, then run npm run generate:books.
% canonical-source-hash: fnv1a64:f3e41532470eb9fb
% canonical-lessons: IT-C343-dubitare, IT-C343-scettico, IT-C343-negare, IT-C343-assurdo, IT-C343-discutibile

\chapter{Ho i miei dubbi}
\label{ch:it-ho-i-miei-dubbi}
\hlchaptermodality{343}

\begin{chapteropening}
\textbf{By the end of this chapter:} \emph{I can say I doubt something, that I am sceptical, that an idea is absurd or questionable, and deny a claim.}

Complete the last lesson of chapter 343: I can say I doubt something, that I am sceptical, that an idea is absurd or questionable, and deny a claim.
\end{chapteropening}

\section[dubitare]{dubitare — to doubt}

\label{lesson:IT-C343-dubitare}



\begin{quote}
Before the new one: say the Italian for the compromise, then the Italian for to give in.
\end{quote}

\subsection*{You'll want to know: dubitare}
\textbf{dubitare} — \textquotedblleft{}to doubt\textquotedblright{}.

\begin{quote}
Dubito che arrivi in orario.
\end{quote}

\begin{quote}
\emph{I doubt that he will arrive on time.}
\end{quote}

\emph{Dubitare che} takes the subjunctive. \emph{Dubitare di qualcuno} — not to trust someone.

\begin{cousinweb}
\emph{dubitare} is a learned borrowing from Latin \emph{dubitāre} \textquotedblleft{}to waver, to be uncertain, to doubt\textquotedblright{} (EtymDB), related to \emph{dubius} \textquotedblleft{}doubtful\textquotedblright{}, a word linked with \emph{duo} \textquotedblleft{}two\textquotedblright{}: a mind pulled two ways.
\end{cousinweb}

\begin{grammarlens}[title={What you've built}]
One more word for telling, explaining and describing at B1.
\end{grammarlens}

\subsection*{Your turn}
\begin{itemize}
\raggedright
  \item \emph{Say it:} \emph{dubitare}
  \item \emph{Say it:} \emph{dubitare}, once more
  \item \emph{Say it:} a sentence of your own with \emph{dubitare}
  \item \emph{Recall:} say the Italian for the compromise, then the Italian for to give in, then say \emph{dubitare} again
\end{itemize}

\begin{tcolorbox}[breakable,colback=teal!4,colframe=teal!35!black,title={Before you move on}]
What does dubitare mean? (\textquotedblleft{}To doubt\textquotedblright{}.) Say it once more.
\end{tcolorbox}
\section[scettico]{scettico — sceptical}

\label{lesson:IT-C343-scettico}



\begin{quote}
Before the new one: say the Italian for to give in, then the Italian for to doubt.
\end{quote}

\subsection*{You'll want to know: scettico}
\textbf{scettico} — \textquotedblleft{}sceptical\textquotedblright{}.

\begin{quote}
Sono un po' scettico su questa dieta.
\end{quote}

\begin{quote}
\emph{I am a bit sceptical about this diet.}
\end{quote}

Plural \emph{scettici}, feminine \emph{scettica}. \emph{Essere scettico su qualcosa}.

\begin{cousinweb}
\emph{scettico} is a learned borrowing, through Latin \emph{scepticus}, from Greek \emph{skeptikós} \textquotedblleft{}inquiring, reflective\textquotedblright{}, from \emph{sképtesthai} \textquotedblleft{}to look about carefully\textquotedblright{}; a related Greek word passed into Latin as \emph{scopos} \textquotedblleft{}a mark to aim at\textquotedblright{} (Lewis \& Short).
\end{cousinweb}

\begin{grammarlens}[title={What you've built}]
One more word for telling, explaining and describing at B1.
\end{grammarlens}

\subsection*{Your turn}
\begin{itemize}
\raggedright
  \item \emph{Say it:} \emph{scettico}
  \item \emph{Say it:} \emph{scettico}, once more
  \item \emph{Say it:} a sentence of your own with \emph{scettico}
  \item \emph{Recall:} say the Italian for to give in, then the Italian for to doubt, then say \emph{scettico} again
\end{itemize}

\begin{tcolorbox}[breakable,colback=teal!4,colframe=teal!35!black,title={Before you move on}]
What does scettico mean? (\textquotedblleft{}Sceptical\textquotedblright{}.) Say it once more.
\end{tcolorbox}
\section[negare]{negare — to deny}

\label{lesson:IT-C343-negare}



\begin{quote}
Before the new one: say the Italian for to doubt, then the Italian for sceptical.
\end{quote}

\subsection*{You'll want to know: negare}
\textbf{negare} — \textquotedblleft{}to deny\textquotedblright{}.

\begin{quote}
Il ragazzo nega di aver rotto il vetro.
\end{quote}

\begin{quote}
\emph{The boy denies having broken the glass.}
\end{quote}

\emph{Negare di aver fatto qualcosa} — to deny having done something.

\begin{cousinweb}
\emph{negare} is inherited from Latin \emph{negāre} \textquotedblleft{}to say no, to deny, to refuse\textquotedblright{} (Lewis \& Short), a verb built on the same \emph{ne-} \textquotedblleft{}not\textquotedblright{} as \emph{nōn}.
\end{cousinweb}

\begin{grammarlens}[title={What you've built}]
One more word for telling, explaining and describing at B1.
\end{grammarlens}

\subsection*{Your turn}
\begin{itemize}
\raggedright
  \item \emph{Say it:} \emph{negare}
  \item \emph{Say it:} \emph{negare}, once more
  \item \emph{Say it:} a sentence of your own with \emph{negare}
  \item \emph{Recall:} say the Italian for to doubt, then the Italian for sceptical, then say \emph{negare} again
\end{itemize}

\begin{tcolorbox}[breakable,colback=teal!4,colframe=teal!35!black,title={Before you move on}]
What does negare mean? (\textquotedblleft{}To deny\textquotedblright{}.) Say it once more.
\end{tcolorbox}
\section[assurdo]{assurdo — absurd}

\label{lesson:IT-C343-assurdo}



\begin{quote}
Before the new one: say the Italian for sceptical, then the Italian for to deny.
\end{quote}

\subsection*{You'll want to know: assurdo}
\textbf{assurdo} — \textquotedblleft{}absurd\textquotedblright{}.

\begin{quote}
È assurdo pagare così tanto per un caffè.
\end{quote}

\begin{quote}
\emph{It is absurd to pay so much for a coffee.}
\end{quote}

\emph{Ma è assurdo!} — that is ridiculous! Feminine \emph{assurda}.

\begin{cousinweb}
\emph{assurdo} is a learned borrowing from Latin \emph{absurdus} \textquotedblleft{}out of tune, incongruous\textquotedblright{} (EtymDB), built from \emph{ab} \textquotedblleft{}away from\textquotedblright{} and \emph{surdus} \textquotedblleft{}deaf, dull-sounding\textquotedblright{}: a note that does not fit the tune.
\end{cousinweb}

\begin{grammarlens}[title={What you've built}]
One more word for telling, explaining and describing at B1.
\end{grammarlens}

\subsection*{Your turn}
\begin{itemize}
\raggedright
  \item \emph{Say it:} \emph{assurdo}
  \item \emph{Say it:} \emph{assurdo}, once more
  \item \emph{Say it:} a sentence of your own with \emph{assurdo}
  \item \emph{Recall:} say the Italian for sceptical, then the Italian for to deny, then say \emph{assurdo} again
\end{itemize}

\begin{tcolorbox}[breakable,colback=teal!4,colframe=teal!35!black,title={Before you move on}]
What does assurdo mean? (\textquotedblleft{}Absurd\textquotedblright{}.) Say it once more.
\end{tcolorbox}
\section[discutibile]{discutibile — questionable, debatable}

\label{lesson:IT-C343-discutibile}



\begin{quote}
Before the new one: say the Italian for to deny, then the Italian for absurd.
\end{quote}

\subsection*{You'll want to know: discutibile}
\textbf{discutibile} — \textquotedblleft{}questionable, debatable\textquotedblright{}.

\begin{quote}
La scelta dell'allenatore è discutibile.
\end{quote}

\begin{quote}
\emph{The coach's choice is questionable.}
\end{quote}

It politely suggests something is wrong: \emph{una scelta discutibile}.

\begin{cousinweb}
\emph{discutibile} is built on \emph{discutere} \textquotedblleft{}to discuss\textquotedblright{}, from Latin \emph{discutere} \textquotedblleft{}to strike apart, to break up, to scatter\textquotedblright{} (Lewis \& Short), built from \emph{dis-} and \emph{quatere} \textquotedblleft{}to shake\textquotedblright{}: a matter that can be shaken.
\end{cousinweb}

\begin{grammarlens}[title={What you've built}]
One more word for telling, explaining and describing at B1.
\end{grammarlens}

\subsection*{Your turn}
\begin{itemize}
\raggedright
  \item \emph{Say it:} \emph{discutibile}
  \item \emph{Say it:} \emph{discutibile}, once more
  \item \emph{Say it:} a sentence of your own with \emph{discutibile}
  \item \emph{Recall:} say the Italian for to deny, then the Italian for absurd, then say \emph{discutibile} again
\end{itemize}

\begin{tcolorbox}[breakable,colback=teal!4,colframe=teal!35!black,title={Before you move on}]
What does discutibile mean? (\textquotedblleft{}Questionable, debatable\textquotedblright{}.) Say it once more.
\end{tcolorbox}
